\subsubsection{Quá trình đẳng tích}
\textit{Quá trình đẳng tích} là quá trình biến đổi trạng thái của chất khí trong đó thể tích V được giữ nguyên không đổi.
\\
\textit{Trong quá trình đẳng tích của một lượng khí nhất định, áp suất tỉ lệ thuận với nhiệt độ tuyệt đối}
\begin{align*}^{}
	p\sim T\ hay\ \dfrac{p}{T}=\ct{ hằng số}
\end{align*}
\par Khi chất khí chuyển từ trạng thái này sang trạng thái khác thì ta có:
\begin{align*}\boder{\dfrac{p_1}{T_1}=\dfrac{p_2}{T_2}= \dfrac{p_3}{T_3}=..}\end{align*}
\subsubsection{Đường đẳng tích}
\begin{itemize}
	\item \textbf{Định nghĩa:} Là đường biểu diễn sự phụ thuộc của áp suất theo nhiệt độ tuyệt đối khi thể tích không đổi.
	\item Dạng đường đẳng tích trong các hệ trục toạ độ:
	\begin{center}
		\begin{tikzpicture}
			\draw[->] (0,0)node[left]{O}--(3,0)node[below]{$T$}; 
			\draw[->] (0,0)--(0,3)node[left]{p};
			\draw[dashed] (0,0)--++(50:1);
			\draw [blue, thick](0,0)++(50:1)--++(50:2)node[above,black]{$V_1$};
			\draw[dashed] (0,0)--++(30:1);
			\draw[blue, thick] (0,0)++(30:1)--++(30:2)node[right,black]{$V_2$};
			\draw (0,0)++(40:3)node[above right]{$V_1<V_2$};
			\node at (1.5,-1){Hệ trục toạ độ (p, T)};
			\begin{scope}[shift={(5,0)}]
				\draw[->] (0,0)--(3,0)node[below]{$V$}; 
				\draw[->] (0,0)--(0,3)node[left]{$p$};
				\draw [blue, thick] (2,1)--++(90:2);
				\node at (1.5,-1){Hệ trục toạ độ (p, V)};
			\end{scope}
			\begin{scope}[shift={(10,0)}]
				\draw[->] (0,0)--(3,0)node[below]{$V$};
				\draw[->] (0,0)--(0,3)node[left]{$T$};
				\draw [blue, thick] (2,1)--++(90:2);
				\node at (1.5,-1){Hệ trục toạ độ (T, V)};
			\end{scope}
		\end{tikzpicture}
	\end{center}
\end{itemize}